% spellchecker:disable
\documentclass[pdftex,
	chapterprefix,
	headsepline,
	footsepline,
	% colordvi,
	11pt,
	a4paper,
	parskip=half,
	% enabledeprecatedfontcommands,
	final,
	appendixprefix,
	bibliography=totoc]{scrbook}

% define a new if command to distinguish between PDF and DVI output
\usepackage{ifpdf}
\ifx\pdfoutput\undefined
\pdffalse %not PDFLaTeX
\else
\pdfoutput=1
\pdftrue
\fi

% language support (hyphenation etc)
% this sets the main language of the paper to english
% if you wish to compose it in german, switch the options around
\usepackage[german,english]{babel}

% for prettier tables
\usepackage{booktabs}

% support for latin1 characters. That means you can enter umlauts directly
% no need for "a "u "o "s anymore
\usepackage[utf8]{inputenc}
\usepackage{textcomp}

% provides the \url{} command to pretty print urls
\usepackage{url}

% allows to \includegraphics
\usepackage{graphicx}
\usepackage{svg}

% defines some standard colornames like "black" etc.
\usepackage{color}

% allows to color tablecells
\usepackage{colortbl}

% provides an easier interface to if-then-else constructs in custom macros
\usepackage{ifthen}

% allows tables to break over pages.
\usepackage{supertabular}

% allows to have different kinds paper orientations in the same pdf-documnent
\usepackage{pdflscape}

% allows to specify absolute texpos for textboxes. This is generally only important for the titlepage
\usepackage[absolute]{textpos}

% % allows to enumerate different figures with a) b) in the same figure-environment.
\usepackage{subfigure}

\usepackage{caption}
\DeclareCaptionLabelFormat{cont}{#1~#2\alph{ContinuedFloat}}
\captionsetup[ContinuedFloat]{labelformat=cont}

% math
\usepackage{amsmath}

% allows inserting source code
\usepackage{listings}
\usepackage{scrhack} % fixes lst incompatibility
\lstset{
% numbers=left,              % Location of line numbers
	numberstyle=\tiny,          % Style of line numbers
	numbersep=5pt,              % Margin between line numbers and text
	tabsize=4,                  % Size of tabs
	extendedchars=true,
	breaklines=true,            % Lines will be wrapped
	showspaces=false,
	showtabs=false,
	showstringspaces=false
}

% More fancy syntax highlighting: https://texdoc.net/texmf-dist/doc/latex/minted/minted.pdf
% careful, this has additional external dependencies
%\usepackage[newfloat]{minted}
% configure source code highlighting; edit this to your liking
% style=bw will save you some money on color pages, but may be harder to read
%\setminted{
%	style=stata-light,
%	frame=leftline,
%	breaklines=true,
%	numbers=left,
%	tabsize=4}

% allows usage of BibLaTeX for bibliography; needs the Biber tool!
\usepackage[
backend=biber,
sorting=anyt
]{biblatex}
\usepackage{csquotes}
% point this to your bibliography file
\addbibresource{bib/references.bib}

% finetune the gaps between figure and text in the subfigure environment (basically close the gap as much as possible)
\renewcommand{\subfigtopskip}{0pt}
\renewcommand{\subfigbottomskip}{0pt}

% some color definitions for the pdf statements below
\definecolor{mygrey}{rgb}{0.45,0.45,0.45}
\definecolor{mydarkgrey}{rgb}{0.2,0.2,0.2}
\definecolor{red}{rgb}{1.0,0.33,0.33}
\definecolor{orange}{rgb}{1.00,0.73,0.33}
\definecolor{yellow}{rgb}{0.95,0.92,0.}
\definecolor{lightgreen}{rgb}{0.3,0.95,0.46}
\definecolor{titleblue}{rgb}{0.03,0.10,0.46}

\ifpdf
% For screen viewing it is nice to have references marked in a slightly different
% color than the rest of the text. Since they will be hyperlinks to the
% referenced objects.
\usepackage[pdftex,
	pdftitle={Fancy thesis template},
	% colorlinks,
	% linkcolor={black},
	% citecolor={black},
	urlcolor={black},
	plainpages={false},
	bookmarksnumbered={true},
	pdfauthor={Vorname Nachname},
	pdfsubject={},
	pdfkeywords={},
	pdfstartview={FitBH}]{hyperref}

\pdfcompresslevel=9
\fi

% some configuration for the amount of text on a single page
\usepackage{typearea}
\areaset[1.5cm]{418pt}{658pt}
\setlength{\headheight}{37pt}

% enter author and title for the titlepage.
\author{}
\title{}

% To avoid nasty mistakes like having comments directly in the textflow
% the following \todo macro was defined. With that you can enter
% \todo{What I still have to do here}
% inside of your text and a marker will appear at the page's margin with the
% text "What I still have to do here".
% The first line activates this feature. If you comment it out and uncomment
% the second line below there will be no error messages and no todos will be shown
% anymore. So - even if you have forgotten to delete one of them - they will not appear
% in the final printout.
\newcommand{\todo}[1]{\marginpar{\textcolor{red}{TODO:} #1}}

% We recommend to split your document into several files. Usually one for every chapter is a
% good idea. If you follow this guideline (how to assemble these files in a single document
% see two paragraphs below) you will be able to single out chapters via the \includeonly{}
% command. Using this mechanism page numbering and references of the full run before will be
% preserved. This also nice, if your latex run tends to get slow and you need to fine tune
% some formatting in one chapter - just include that one. The rest (or at least the ones before
% the one currently under construction) will remain untouched. This means a boost in compilation time.
%\includeonly{content/chapter2}

\begin{document}

% the next two lines influence the detailedness of the table of contents
% and to what structure depth numbers are written before sections/subsections/paragraphs
% You should not touch this
\setcounter{tocdepth}{2}
\setcounter{secnumdepth}{3}

\frontmatter

% here the titlepage is included. Look into the file "frontpage.tex" to
% adapt it to your needs (name, title etc.)
% spellchecker:disable
\begin{titlepage}
\selectlanguage{german}
\vspace*{-1cm}
\newlength{\links}
\setlength{\links}{0.9cm}
\setlength{\TPHorizModule}{1cm}
\setlength{\TPVertModule}{1cm}
\textblockorigin{0pt}{0pt}

\sffamily
\LARGE

\begin{textblock}{16.5}(2.8,2.6)
 \hspace*{-0.25cm} \textbf{UNIVERSITÄT DUISBURG-ESSEN} \\
 \hspace*{-1.15cm} \rule{5mm}{5mm} \hspace*{0.05cm} FAKULTÄT FÜR INGENIEURWISSENSCHAFTEN\\
 \large{}ABTEILUNG INFORMATIK UND ANGEWANDTE KOGNITIONSWISSENSCHAFT\\
\end{textblock}


%Hier Titel, Name, und Matrikelnummer eintragen, \\ make a newline
\begin{textblock}{14.5}(3.2,8.5)
  \large
{ \textbf{Bachelorarbeit}} \\[1cm]
{\LARGE \Large\textbf{An Algorithm for Dependency-Preserving Smart Home Updates}} \\[1.3cm]
Koray Uzun\\
Matrikelnummer:3081690\\
Angewandte Informatik (Bachelor)
\end{textblock}



\begin{textblock}{10}(10.5,17.5)
\includegraphics[scale=1.0]{images/unilogo.pdf}\\
\normalsize
\raggedleft
Fachgebiet Verteilte Systeme, Abteilung Informatik \\
Fakultät für Ingenieurwissenschaften \\
Universität Duisburg-Essen \\[2ex]

\today\\[15ex]
\raggedright
% Supervisors
{\textbf{Erstgutachter:}} Prof. Dr-Ing. Torben Weis \\
{\textbf{Zweitgutachter:}} Peter Zdankin \\
{\textbf{Zeitraum:}} 1. September 2042 - 1. Januar 2043 \\
\end{textblock}

\end{titlepage}


\selectlanguage{english}

\chapter*{Abstract}


Smart Home Systeme sind Netzwerke von elektronischen Geräten innerhalb eines Wohnraumes. 
Diese Systeme unterliegen, wie jede andere Technologie auch, einem ständigen Wandel. Während andere Technologien,
wie zum Beispiel Smartphones meist nach einigen wenigen Jahren ausgetauscht, erwartet man von
Smart Home Systemen eine lange Lebensdauer.
Die meist komplexen und teuren Installationen sind nicht dazu ausgelegt alle paar Jahre 
ausgetauscht zu werden. Dennoch erwartet man, wie von anderen Technologien auch, dass die Systeme
sich den Gegebenheiten der Zeit anpassen. Neue Funktionen oder oft auch Sicherheitslücken erfordern 
ein ständiges updaten der Systeme oder einzelner Geräte. Mit diesen Updates geht jedoch ein großes Risiko 
einher. Aufgrund der riesen Diversität verschiedener Smart Home Systeme kann vom Hersteller meist nicht
garantiert werden, dass ein Update keine Schäden anrichtet. Lediglich das Update im Einzelnen wird auf 
mögliche Fehlerquellen überprüft, nicht jedoch das Zusammenspiel des Updates mit anderen Geräten.
Dies wäre aufgrund der riesen Diversität an Smart Home Systemen für den Hersteller auch gar nicht umsetzbar.
Dennoch benötig man eine Möglichkeit Updates präventiv zu überprüfen, um so größere Schäden zu meiden und 
die Langlebigkeit von Smart Home Systemen zu gewährleisten.
Diese Abschlussarbeit beschäftigt sich mit der Implementation und Evaluation eines Algrorithmus, welcher 
Abhängigkeiten zwischen Geräten untersucht und basierend auf diesen Abhängigkeiten Updatekonfigurationen 
findet, die keine Probleme verursachen.
Der Algorithmus basiert auf dem Paper "An Algorithm for Dependency-Preserving 
Smart Home Updates" von Peter Zdanking, Matthias Schaffeld, Marian Waltereit, Oskar Carl und Torben Weis. 






\footnote{Wikipedia: \url{https://en.wikipedia.org/wiki/Abstract_(summary)}}



\tableofcontents

%\listoffigures
\mainmatter

% To assemble the whole document
% Please be aware that each file will begin on a new page
% therefore chapters should be put into such a file.
% There cannot be an include statement inside of an "included" file.
% So if you want to further divide your document - use \input inside of
% the included files. \input will not begin on a new page.
\chapter{Einführung}

In den letzten Jahren stieg die Anzahl der Smart Home Geräte stark. Laut Verified Market Research wird
der gloable Smart Home Markt im Jahr 2019 auf bereits 80 Milliarden Dollar geschätzt. Prognostiziert werden 207.88 
Milliarden Dollar im Jahr 2027. Dies entspricht einer jährlichen Wachstumsrate von 13.52 Prozent [2].
Zudem ist es so, dass innerhalb von Smart Home Systemen die Anzahl der Geräte ebenfalls stetig steigt. Während 2016
durchschnittlich 5.8 Geräte über ein Smart Home System vernetzt waren, waren es 2018 schon bereits 8.1 Geräte [3]. 
Es ist zu vermuten, dass sich dieser Trend vorerst so weiterentwickeln wird und somit Smart Home Systeme immer
größer werden. Wie jede andere Technologie müssen Smart Home Systeme regelmäßig geupdatet werden, um zum Beispiel
Sicherheitslücken zu schließen oder neue Funktionalitäten hinzuzufügen. 
Mit der steigenden Anzahl an Geräten, steigt in einem Smart Home System auch logischerweise die Anzahl der Updates und somit 
insgesamt die Komplexität von Smart Home Systemen. Das Problem hierbei ist, dass mit steigender Komplexität zum einen automatisch eine höhere
Fehleranfälligkeit bezüglich Updates einhergeht, zum anderen wird es schwieriger Fehler zu beheben,
die durch Updates verursacht werden. Updates im Einzelnen besitzen bereits viele Fehlerquellen, wie zum Beispiel Sicherheitslücken 
oder generell fehlerhafte Implementationen. Dies sind jedoch Aspekte die bereits vom Entwickler weitestgehend vermieden werden.
Viel Problematischer ist das Zusammspiel verschiedener Geräte und den dazugehörigen Updates. Ein an sich fehlerfreies Update
kann bei der Installtion in einem Smart Home System, trotz "fehlerfreier" Programmierung große Probleme verursachen. Was ist wenn zum Beispiel 
ein Update eine Funktion eines Geräts so verändert, dass ein anderes Gerät nicht mehr darauf zugreifen kann. Durch solche kleine Änderungen
können Abhängigkeiten zwischen Geräten zerstört werden, wodurch ganze Teile eines Smart Home Systems ausfallen
können. Solch ein Ausfall bedeutet für den Nutzer meist ein großer finanzieller Schaden, da Updates oft nicht 
rückgängig gemacht werden können und somit Ersatz beschafft werden muss.
Aus diesen Gründen sind präventive Maßnahmen notwendig,  um die Langlebigkeit von Smart Home Systemen gewährleisten zu können.
Langlebigkeit wird nämlich von den meisten Usern erwartet, wie man an der Resonanz zum Abschalten der HueBridge V1 von Philips sehen kann 
[4].


\section{Ziel der Arbeit}


Eine Steigerung der Lebensqualität ist das Kernziel von Smart Home Systemen. Dies erreicht man zum Beispiel durch die Konfiguration von 
automatisierten Prozessen. Dadurch können Nutzer nämlich Zeit, Energie und in gewissen Belangen auch Geld sparen. Laut Umfragen geben 57 Prozent der Nutzer von
Smart Home Systemen in Amerika an jeden Tag durchschnittlich 30 Minuten Zeit einzusparen [4].
Zusätzlich lassen sich nach Berechnungen des Fraunhofer Instituts für Bauphysik  durch Automatisierungen bis zu ca. 40 Prozent Heizkosten 
sparen. [5].
Solche Vorteile zeigen, dass sich die anfangs teure Anschaffung eines Smart Home Systemes auf lange Zeit sogar rentieren kann. Sicherlich 
ist dies der Idealfall, aber dennoch zeigen diese Zahlen, dass Smart Home Systeme viele Potenziale bieten. Es hängt jedoch von 
vielen Faktoren ab, ob sich die Anschaffung eines Smart Homes lohnt und objektiv ist dies nicht zu beurteilen. Was jedoch 
objektiv gesagt werden kann, ist dass in den meisten Fällen Nutzer aufgrund der teuren Anschaffungskosten eines Smart Homes davon ausgehen, dass 
ihr System viele Jahre erhalten bleibt. Zum Beispiel gibt es Smart Home Systeme, welche beim Bau eines Hauses direkt mit geplant werden,
sodass man, wie zum Beispiel bei KNX Systemen alle Geräte über ein permanent installiertes Bus System verbinden kann [1].
Selbstverständlich gibt es auch günstigere Varianten, aber dennoch ist die generelle Erwartungshaltung, dass wenn man ein funktionierendes System hat, dieses
auch über viele Jahre hinweg weiter funktionieren wird.
Daher is das Ziel dieser Arbeit die Implementation und Evaluation eines Algrorithmus, welcher 
Abhängigkeiten zwischen Geräten untersucht und basierend auf diesen Abhängigkeiten Updatekonfigurationen 
ermittelt, welchen den größten Nutzen für den Nutzer haben.
Mit diesem Algorithmus wäre es für Smart Home Nutzer möglich Updates bereits vor der Installation zu überprüfen, um so sicherzustellen, 
dass sie keine Ausfälle verursachen. Die Anwendung dieses Algorithmus sollte für den Nutzer möglichst einfach sein. So könnte man zum 
Beispiel ein Smartphone nutzen, welches als zentrales Gerät im System dient. Auf diesem zentralen Gerät wäre es dann möglich den 
Algorithmus laufen zu lassen und sich über das Smartphone dann für die passende Updatekonfiguration zu entscheiden. Smartphones
bieten sich besonders an, da sie eine intuitive Bedienung für den Nutzer ermöglichen und im Normalfall jeder Nutzer eines Smart Homes
ein Smartphone besitzt. So umgeht man das Problem, dass es, wie bereits erwähnt, etliche Smart Home Anbieter mit verschiedener
Software und Geräten gibt. 




 


\chapter{Related Work}\label{ch:related_work}
In diesem Kapitel wird das Thema Langlebigkeit im Bereich Software Engineering thematisiert. Infolgedessen werden
Smart Home Systeme unter dem Aspekt der Langlebigkeit genauer beleuchtet und  das Thema Dependency Management aufgegriffen, wobei bereits 
bestehende Ansätze/Lösungen diesbezüglich genauer betrachtet werden. Zusätzlich wird der aktuelle Stand des Dependency Managements im 
Bereich Smart Home untersucht.


\section{Smart Homes}
Smart Home Systeme sind grundlegend nichts anderes als die Verknüpfung elektronischer Geräte. Durch die hohe Anzahl 
der Anbieter für Smart Home Systeme existieren jedoch viele Architekturen, die teilweise sehr unterschiedlich funktionieren .
Im Rahmen dieser Arbeit ist es nicht notwendig die verschiedenen Smart Home Architekturen genauer zu beleuchten,
dennoch schadet es nicht sich einen groben Überblick über die aktuellen Technologien zu verschaffen. Grob können die
verschiedenen Architekturen darin unterscheiden werden, ob sie eine Verknüpfung zu einer Cloud besitzen. Cloud gebundene
Architekturen bieten dem Nutzer meist eine einfachere Installation und Wartung, indem zum 
Beispiel Updates automatisch installiert werden. Gleichzeitig bringen sie aber auch potenzielle Gefahren mit sich, da die
Abhängigkeit zu einer Cloud Smart Home Systeme verwundbar machen kann. Wird ein externer Service, wie zum
Beispiel die Cloud abgeschaltet, führt dies bei Cloud-gebundenen Archtiekturen oft zu Ausfällen im System[7].
Des Weiteren sind Geräte, die mit einer Cloud verbunden sind anfälliger gegenüber Angriffen von Außenstehenden, die sich 
Zugriff auf das System beschaffen wollen.
Nicht Cloud-gebundene Architekturen haben diese Nachteile nicht. Es ist jedoch meist schwieriger eine solche Architektur aufzubauen, 
da zum Beispiel, wie im Falle von OpenHab[?] viel manuelle Konfiguration benötigt wird. Zusätzlich bieten nicht Cloud-gebundene Architekturen
mehr Sicherheit, da sie keine starken Abhängigkeiten zu anderen Dienstleistern benötigen.
Natürlich könnte man Smart Home Systeme noch weiter spezifizieren, indem man einzelne Komponenten genauer beleuchet und 
verschiedene Smart Home Anbieter miteinander vergleicht, jedoch ist dies, wie bereits erwähnt, im Rahmen dieser 
Arbeit nicht notwendig. Der zu entwickelnde Algorithmus soll später theoretisch Platform- und Architekturunabhängig anwendbar sein.
Zu diesem Zweck werden im Folgenden einige Aspekte bezüglich Smart Homes definiert und vorausgesetzt. 
Zdankin beschreibt in seiner Arbeit Smart Home Systeme als:

\par
\begingroup
\leftskip=1cm % Parameter anpassen
\noindent %ab hier der Text, der eingerückt werden soll
\textit{"A smart home has a set of devices that are connected through
a platform. Each device has a certain software version and
a set of available updates. Each software version has a set
of available predefined services. Devices can use services of
other devices which creates dependencies. An update configuration
is one of the finite states of the nondeterministic
finite automaton (NFA) that can be constructed by using the
configurations as states and connecting them using individual
updates as transitions." (Zdankin et al., 2020, Requirements and Mechanisms for Smart Home Updates [8])}
\par
\endgroup
Um Missverständnisse zu vermeiden folgen nun zusätzlichen Begriffserklärungen.
Die folgenden Elemente in Abbildung 2.1 sind für ein Smart Home System also essentiell:
\begin{itemize}
\item Device (Gerät): Geräte sind die Komponenten, die die eigentliche Funktionalität anbieten. Sie besitzen eine aktuelle Firmware
Version mit entsprechenden Diensleistungen(Services) und potenzielle Updates. Zusätzlich können sie Abhängigkeiten zu anderen Geräten haben, indem sie auf deren
Dienstleistungen zugreifen.
\item Firmware Version: Die Firmware Version gibt an, auf welcher Version ein Gerät sich befindet und gibt somit an welche Dienstleistungen ein Gerät anbietet.
\item Service (Dienstleistungen): Eine Dienstleistung ist eine Funktionalität, die ein Gerät anbietet. Zum Beispiel bietet ein Thermostat
die Dienstleistung "Temperatur messen" an.
\item Update: Eine Aktualisierung der Firmware eines Geräts. Durch ein Update können sich angebotenen Diensleistungen eines Geräts ändern.
\item Platform: Die Komponente, die alle Geräte miteinander vernetzt.
\item Dependency (Abhängigkeiten): Eine Abhängigkeit zwischen zwei Geräten besteht, wenn ein Gerät Zugriff auf die Dienstleistungen eines
anderen Geräts benötigt.
\end{itemize}

\begin{figure}[!ht]
\hspace*{-2,5cm}                                                           
\includegraphics{"System Model"}
  \caption{Grundlegende Komponenten eines Smart Home Systems}
  \label{fig:Prob1:MEA}
\end{figure}


\newpage
\subsection{Prescriptive vs Descriptive Standards}
Für die Kommunikation der Geräte untereinander existieren viele Kommunikationsprotokolle wie zum Beispiel Bluetooth,
Wi-Fi oder auch XMPP. Diese sind bereits im großen Umfang in Verwendung und stark erforscht, weshalb es keiner weiteren
Untersuchung dieser Protokolle im Rahmend dieser Arbeit bedarf.
Kommunikation allein recht jedoch nicht, die Geräte müssen nicht nur untereinander kommunizieren, sondern sie müssen sich auch 
verstehen. Dafür ist ein Standard notwendig, welcher die Funktionen/Dienstleistungen der Geräte beschreibt.
Die zwei führenden Ansätze im Bereich IoT/Smart Homes sind deskriptive und preskriptive Standards.
Zdankin beschreibt die beiden Standards in seiner Arbeits folgendermaßen:

\textit{
"A standard for service definitions that only regulates how a device can describe itself is
a descriptive standard, because it allows vendors to describe their devices in their own
terms. A recently finalized descriptive standard is the Web Of Things [Ka20; Ko20]. If
devices need to communicate across different descriptive standards, a translation must be
considered. However, McCool et al. have stated security concerns about purely descriptive
approaches, such as scanning for door locks with known physical weaknesses [MR18] (Zdankin et al., 2020, Requirements and Mechanisms for Smart Home Updates [8]).}
\todo{Wie muss ich die Quellenangaben im Zitat zitieren?}

%\footnotetext{[Ka20]: httpsȷ//www.w3.org/TR/wot-thing-description/}
%\footnotetext{[Ko20]: httpsȷ//www.w3.org/TR/wot-architecture/}
%\footnotetext{[MR18]: Mccool, M.; Reshetova, E.ȷ Distributed Security Risks and Opportunities in the
%W3C Web of Things. In. Jan. 2018.}

\textit{
A prescriptive standard does not only regulate the how but also the what of self-description.
These standards prescribe how devices and services should be defined and vendors can use
these terms for unambiguity. By using the same terms to define services, translation is not
required anymore. A prescriptive standard defines all possible devices and all their services
they could implement. However, a device can implement a subset of these services only.
This restriction has benefits as well because it ensures compatibility across device vendors.
Widely used examples are the smart home standards of Amazon, Apple, and Samsung6."
(Zdankin et al., 2020, Requirements and Mechanisms for Smart Home Updates [8])}

In der Realität bieten sich hybride Standards vermutlich am meisten an, da sowohl deskriptive als auch preskriptive Standards Vor- und
Nachteile besitzen.   


\subsection{Update Planing}
Oft werden Smart Home Geräte automatisch über eine aktive Internetverbindung geupdatet oder auch durch andere
Geräte wie zum Beispiel einem Hub. Dies bietet dem Nutzer mehr Komfort, da kein manueller Aufwand für Updates ensteht.
Solches "Remote Updating" sollte aber eigentlich vermieden werden, da dadurch zum Beispiel schädliche Software installiert werden
kann. Außerdem müsste der externe Diensleister, der das Updaten der Geräte übernimmt, viele Informationen über das System sammeln,
um so ideale Updatekonfigurationen zu finden. Das Preisgeben aller Informationen sollte nicht leichtsinnig getan werden, da es viele
potenzielle Gefahren birgt. Allein die Information, wann eine automatisierte Heizung anfängt zu heizen kann potenziellen Kriminellen Informationen
darüber geben, wann jemand zu Hause ist oder wann ein Haus leer steht. 
Idealerweise sollten Nutzer daher das Updaten lokal selbst übernhemen. Dies bedeutet selbstverständlich mehr manuelle Konfiguration,
vor jedem Update, da der Nutzer gefragt werden sollte, ob er ein Update installieren möchte. Diese Frage kann ein Nutzer ohne weitere Informationen
gar nicht plausibel entscheiden. Es wäre also von Vorteil, wenn es für den Nutzer
eine Übersicht über alle möglichen Updates geben würde, in der er sich dann für passende Updates entscheiden kann. 
Für einen solchen Zweck bietet sich Smartphone als zentrales Gerät an um eine Übersicht über alle im 
System vorkommenden Geräte zu erstellen.
Dazu benötigt man eine Möglichkeit die Informationen über die Services und Updates eines Geräts leicht abfragen zu können.
Eine solche Abfrage müsste in der Realität für jedes Gerät individuell angepasst werden. Dies ist jedoch nicht umsetzbar,
weswegen es nötig ist einen einheitlichen Weg zu finden.
Zdankin schlägt vor die Metadaten der Geräte zu diesem Zweck zu nutzen. Geräte besitzen üblicherweise Metadaten, welche Informationen über 
Software Version, Speicherbedarf usw. enthalten. Diesen Metadaten könnte man zusätzlich Informationen über Dienstleistungen und Updates 
hinzufügen, sodass man einen schnellen Überblick über alle Geräte innerhalb eines Netzwerks erhalten kann
(Zdankin et al., 2020, Requirements and Mechanisms for Smart Home Updates [8]). 

\section{Langlebigkeit}
Im Bereich Software Engineering ist Langlebigkeit ein bekanntes Problem. Es wurden bereits Untersuchungen über die durchschnittliche 
Lebensdauer von Software durchgeführt mit dem Ergebnis, dass kleine Projekte im Bereich von unter 100.000 Zeilen Code wird eine Lebensdauer von circa 
6-8 Jahren haben. Dabei ist zu beobachten, dass mit steigender Größe eines Projekts die erwartete Lebensdauer im Durchschnitt
ebenfalls steigt [9]. Dies hängt vermutlich mit den hohen Kosten zusammen, die für größere Projekte aufkommen, weswegen man daran
interessiert ist lange einen Nutzen von der Software zu haben. Gründe für das Altern von Software sind vielseitig. So kann zum Beispiel 
die Wartung von Software einfach zu teuer werden, so dass es sich lohnt neue Software zu entwickeln oder eine Entwicklungsumgebung wird
nicht mehr unterstützt und es lohnt sich nicht die Software zu portieren. Andere Gründe können Sicherheitslücken oder auch fehlerhafte 
Updates sein. Die German Research Foundation hat ein Projekt namens Design for Future ins Leben gerufen, um das Problem der Langlebigkeit
von Software zu untersuchen. Das Projekt beschreibt sein Ziel darin fundamentale neue Ansätze im Bereich langlebiger Software zu finden, um Probleme mit Legacy-Software
oder der Adaption von Software auf neue Plattformen zu lösen und eine kontinuierliche Weiterentwicklung von Software Systemen gewährleisten 
zu können [10].
Im Bereich Smart Home existieren keine solche Programme. Das Problem der Langlebigkeit wird in diesem Kontext unterschätzt oder einfach ignoriert.
Aufgrund der Komplexität der Systeme besitzen sie viele Schwachstellen die zu einem frühzeitigen Ausfall des Systems führen können. Dazu gehören:

\begin{figure}[h]
\centering
\includegraphics{"Threats"}
\caption{Gefahren für die Langlebigkeit von Smart Home Systemen\( ^1\) }
\label{fig:Prob1:MEA}
\end{figure}
\footnote{Zdankin et al, 2020, An Algorithm for Dependency-Preserving Smart Home Updates}

In dieser Arbeit wird hauptsächlich die Gefahr der "Breaking Updates" thematisiert.
"Breaking Updates" können Abhängigkeiten im System zerstören und sind somit Schäden gegen die man als Nutzer nichts unternehmen
kann. Eine Möglichkeit wäre es Updates generell nicht zu installieren, wodurch Nutzer auf neue Funktionalitäten verzichten müssten, während gleichzeitig das Risiko steigt,
dass Sicherheitslücken im System aufkommen. Updates also grundlegend zu ignorieren ist nicht empfehlenswert, weswegen es einer Möglichkeit bedarf
Updates vor ihrer Installation zu überprüfen.
Eine andere Gefahr ist zum Beispiel, dass ein externer Service, der für das Smart Home System notwendig ist, offline gehen kann (Discontinued External Service).[1]
Der Unteschied zu den "Breaking Updates" besteht darin, dass man als Käufer eines Cloud-gebundenen Smart Home Systems diese Gefahr aktiv und willentlich eingeht.
Die ist bei "Breaking Updates" nicht der Fall.
Außerdem lässt sich das Ausmaß der Gefahr eines "discontinued Services" weitesgehend vermeiden, indem man bei etablierten Anbietern einkauft oder
komplett auf die Abhängigkeit zu einem externen Service verzichtet. Auch dies ist bei "Breaking Updates" nicht möglich.


\newpage
\section{Dependency Managment}
Unter Dependency Management versteht man im Allgemeinen das Strukturieren von Abhängigkeiten verschiedener 
Systeme/Programme. Oft können Programme nur in Abhängigkeit von anderen Programmen starten.
Dies stellt kein Problem dar, solange die Abhängigkeiten in ihrer Anzahl überschaubar bleiben. 
Realität ist jedoch, dass Programme meist von mehreren Programmen abhängig sind und diese
Prorgramme wiederum abhängig von anderen Programmen. Bei zu vielen Abhängigkeiten verliert man 
als Entwickler und auch als Nutzer schnell den Überblick. Dieses Wirrwarr an Abhängigkeiten bezeichnet man umgangssprachlich 
als "Dependency Hell", denn sobald man einmal den Überblick verloren hat, wird es äußerst schwierig weitere Änderungen am
System vorzunehmen.

\subsection{Probleme}
Es folgt eine kleine Zusammenfassung, der am häufigsten auftretenden Formen der "Dependency Hell" [11]:
\begin{itemize}
 \item Long Chain of Dependencies: Ein Programm A ist abhängig von einem Programm B. Das Programm B wiederum ist abhängig von einem Programm C.
Möchte man nun Programm  A nutzen benötigt man zusätzlich Programm  B und damit auch Programm C. Solche lange Ketten können Konflikte verursachen.
(Siehe Conflicting Dependencies)
\item Conflicting Dependencies: Sei ein Programm A abhängig von einem Programm B1.1, ein anderes Programm B ist abhängig von  Programm B1.2. Programm B 
kann aber nicht gleichzeitig mit der Version 1.1 und der Version 1.2 installiert sein. Dies bedeutet, Programm A und B können nicht gleichzeitig in einem System laufen.
\item Circular Dependencies: Sei ein Programm A abhängig von einer spezifischen Version von Programm B. Programm B wiederum ist abhängig von einer spezifischen
Version von Programm A. Beide Programme können nun nicht geupdatet werden, da sie die Abhängigkeiten kaputt verletzen würden.
\item Version Lock: Durch einen Version Lock ist man dazu gewungen beim updaten einer Software alle abhängigen Programme ebenfalls upzudaten. (semver.org)
\item Version Promiscuity: Werden Abhängigkeiten zu locker betrachtet können zukünftige Updates als kompatibel angesehen werden, obwohl sie es nicht sind.
\end{itemize}

Für die genannten Probleme existieren bereits viele Ansätze, um gegen sie vorzugehen oder sie zu vermeiden. Dazu gehören [11]:

\subsection{Lösungen}
\begin{itemize}
 \item Semantic Versioning [12]: Semantic Versioning versioniert Software nach dem Format "MAJOR.MINOR.PATCH", um Probleme wie Version Lock 
oder Version Promiscuity zu vermeiden. Dafür benötigt 
man eine Schnittstelle der man die Änderungen an der Software mitteilt. Bei Änderungen, die die API betreffen und nicht abwärtskompatibel
sind muss die MAJOR Nummer erhöht werden. Bei Änderungen, die abwärtskompatibel sind, muss die MINOR Nummer erhöht werden. 
Änderungen, die die API nicht betreffen erhöhen die Patch Nummer. Es gibt noch viele weitere Spezifikationen bezüglich Semantic
Versioning, das grobe Konzept sollte jedoch verstanden sein. 
\item Package Manager: Package Manager sind Tools die sich um das installieren, updaten und konfigurieren von Programmen automatisch
kümmern, wodurch eine manuelle Konfiguration nicht mehr nötig ist.
\item Portable Applications: Portable Applikationen umgehen die Probleme, die
Abhängigkeiten mit sich bringen. Sie funktionieren selbstständig,
indem alle nötigen Komponenten bereits vorhanden sind.
\end{itemize}

Die Auflistungen sollen zeigen, dass Dependency Management ein allgegenwärtiges Thema ist, welches in den meisten Bereichen bereits
viel Aufmerksamkeit bekommt. Dependency Management umfasst natürlich noch viele weitere Probleme und dazugehörige Lösungen und Lösungsansätze,
im Bereich Smart Home exisitiert jedoch nahezu keine Forschungsgrundlage. Als Nutzer kann man nur hoffen, dass neue Updates keine Schäden anrichten, 
da es keine Möglichkeit gibt sich davor zu schützen. 


\section{Fazit}

Im Bereich Software Engineering wird Langlebigkeit viel Aufmerksamkeit geschenkt. Smart Home Systeme sind selbstverständlich Teil des
Software Engineerings, aber schaut man sich den aktuellen Stand der Forschung der Langlebigkeit explizit im Bereich der Smart Home
Systeme an, wird deutlich das diesbezüglich kaum Forschungsgrundlagen bestehen.
Dabei werden Smart Home Systeme als Luxusgüter angesehen von denen man aufgrund der hohen Preise eine lange Lebensspanne erwartet.
Um dies zu gewährleisten werden generelle Ansätze des Software Engineerings bezüglich Langlebigkeit genutzt, ausreichend ist dies jedoch nicht.
Ohne vernünftiges Dependency Management kann nicht gewährleistet werden, dass Updates keine Ausfälle in Smart Home Systemen verursachen.
Dies sieht man am Beispiel des Logitech Harmony Hub's [13]. Logitech hatte eine \todo{umschreiben}
neue Firmware Version für das Hub veröffentlicht, welche Sicherheitslücken schließen und Bugs fixen sollte.
Nachdem Update kam es vermehrt zu Ausfällen zwischen dem Hub und third-party Geräten. Den Nutzern waren die Hände 
gebunden, sie konnten nur darauf hoffen, dass Logitech die Probleme wieder behebt. Dies hat Logitech in diesem Fall auch getan,
dennoch zeigt dieses kleine Beispiel, wie leichtsinnig Updates von Nutzern installiert werden. Eine andere Möglichkeit besitzt man auch als 
Nutzer eines Smart Homes nicht wirklich, da es wie eingangs erwähnt, zur Zeit keine Möglichkeiten gibt Updates vor der Installation zu 
überprüfen.








\chapter{Evaluation}\label{ch:eval}

The evaluation usually consists of three main steps:
first it defines the goals intended to be achieved by the software in detail;
next is a description of the methodology
used to measure the satisfaction of the software in relation to these goals;
finally, the measurements are depicted and assessed.


\chapter{Conclusion}

The conclusion quickly summarizes the results of the paper
in relation to the previously defined goals and hypotheses.
It usually also includes some information on next steps and further research
that might be required or possible on the presented subject.




% Appendix chapters to be put here. They will be enumerated with capital letters
% if you  did not change the \documentclass options.
\appendix
% \include{content/appendix}

%Bibliography
% We strongly recommend to use bibtex to manage your bibliography. It helps you
% structure your references and helps avoiding missing important data for a correct
% quotation. If you have no other idea jabref (http://jabref.sourceforge.net/)
% might be a good idea (Java runtime environment needed).
\printbibliography

% Include the eidesstattliche Versicherung
% spellchecker:disable
%pagenumbering{null}

\

\cleardoublepage

\


\pagestyle{empty}

\selectlanguage{german}
\textbf{Versicherung an Eides Statt}\\

Ich versichere an Eides statt durch meine untenstehende Unterschrift,
\begin{itemize}
\item[-] dass ich die vorliegende Arbeit - mit Ausnahme der Anleitung durch die Betreuer - selbstständig ohne fremde Hilfe angefertigt habe und
\item[-] dass ich alle Stellen, die wörtlich oder annähernd wörtlich aus fremden Quellen entnommen sind, entsprechend als Zitate gekennzeichnet habe und
\item[-] dass ich ausschließlich die angegebenen Quellen (Literatur, Internetseiten, sonstige Hilfsmittel) verwendet habe und
\item[-] dass ich alle entsprechenden Angaben nach bestem Wissen und Gewissen vorgenommen habe, dass sie der Wahrheit entsprechen und dass ich nichts verschwiegen habe.
\end{itemize}
Mir ist bekannt, dass eine falsche Versicherung an Eides Statt nach \S 156 und \S 163 Abs. 1 des Strafgesetzbuches mit Freiheitsstrafe oder Geldstrafe bestraft wird.
\vfill
Duisburg, \today\\
$\overline{\parbox{4.8cm}{(Ort, Datum)}} ~~~~~~~~~~~~~~~~~~~~~~~~~~~ \overline{\parbox{7cm}{(Vorname Nachname)}}$


\end{document}
